\chapter{Proprietà degli elementi}

Le configurazioni elettroniche ricavate nel capitolo precedente non hanno soltanto valore descrittivo: da esse dipendono in modo diretto le proprietà fisiche e chimiche degli elementi. 
Poiché le configurazioni si ripetono con regolarità al crescere del numero atomico, anche tali proprietà mostrano un andamento periodico, ed è questa la ragione profonda dell'ordinamento della tavola periodica.

Tutti gli andamenti che seguono si spiegano mediante due sole grandezze già introdotte: la carica nucleare efficace $Z_{\text{eff}}$, che aumenta procedendo verso destra lungo un periodo, e il numero quantico principale $n$ degli elettroni di valenza, che aumenta scendendo lungo un gruppo.

\section{Raggio atomico}

Poiché la funzione d'onda non si annulla a distanza finita, l'atomo non possiede un confine definito e il raggio atomico\mkeyword{raggio atomico} non è una grandezza intrinseca, bensì una definizione operativa: si assume come raggio la metà della distanza minima alla quale due atomi dello stesso elemento possono avvicinarsi. 
Il valore dipende quindi dal tipo di legame considerato, ed è per questo che si distinguono un raggio covalente, uno metallico e uno di van der Waals. Per l'alluminio metallico si ha $r = 143 \ pm$, mentre per il cloro legato covalentemente si ha $r = 100 \ pm$.

Il raggio atomico cresce scendendo lungo un gruppo e procedendo verso sinistra lungo un periodo. Le due tendenze hanno cause distinte.

Scendendo lungo un gruppo gli elettroni di valenza occupano orbitali con $n$ crescente, i quali si estendono progressivamente più lontano dal nucleo; gli elettroni interni, che aumentano di numero, schermano efficacemente la carica nucleare, sicché l'aumento di $Z$ non riesce a compensare l'aumento di $n$.

Procedendo verso destra lungo un periodo, al contrario, gli elettroni vengono aggiunti al medesimo guscio e si schermano fra loro in modo assai poco efficace, mentre la carica nucleare aumenta di un'unità per ogni elemento. 
Ne risulta un aumento di $Z_{\text{eff}}$ che contrae progressivamente l'orbitale: il raggio del litio è maggiore di quello del berillio, benché entrambi abbiano gli elettroni di valenza in $2s$. È un risultato controintuitivo, giacché l'atomo si rimpicciolisce pur acquistando elettroni.

I raggi minori si riscontrano quindi nei gas nobili, che occupano l'estremità destra di ciascun periodo.

\newpage

\section{Raggio ionico}

Uno ione possiede un numero di elettroni diverso da quello dell'atomo neutro, e non è pertanto lecito attribuirgli il raggio atomico. Si definisce allora un raggio ionico\mkeyword{raggio ionico}, ricavato dalle distanze interioniche misurate nei cristalli.

Un catione è sempre più piccolo dell'atomo da cui deriva: la perdita di elettroni riduce la repulsione interelettronica e, se l'elettrone rimosso era l'unico del guscio più esterno, ne elimina uno per intero. Un anione è invece sempre più grande, poiché l'aggiunta di elettroni aumenta la repulsione senza alcuna variazione della carica nucleare. Il sodio passa da $186 \ pm$ a $102 \ pm$ nella forma $Na^+$, mentre il cloro passa da $100 \ pm$ a $181 \ pm$ nella forma $Cl^-$.

\section{Energia di ionizzazione}

Si definisce energia di prima ionizzazione\mkeyword{energia di ionizzazione} l'energia necessaria a rimuovere un elettrone da un atomo neutro isolato allo stato gassoso.

\begin{equation}
    X_{(g)} + E_i \longrightarrow X^+_{(g)} + e^-
    \label{eq:ionizzazione}
\end{equation}

Il processo è sempre endotermico, poiché l'elettrone è legato al nucleo eoccorre fornire energia per allontanarlo indefinitamente.

L'energia di ionizzazione cresce procedendo verso destra lungo un periodo, per effetto dell'aumento di $Z_{\text{eff}}$, e diminuisce scendendo lungo un gruppo, poiché l'elettrone da rimuovere si trova mediamente più lontano dal nucleo. 
I valori massimi si riscontrano quindi nei gas nobili e i minimi nei metalli alcalini: si tratta dello stesso andamento del raggio atomico, percorso in senso inverso.

L'andamento lungo il periodo presenta tuttavia due irregolarità sistematiche, che costituiscono una conferma indipendente della struttura elettronica.
Passando dal gruppo $2A$ al gruppo $3A$ l'energia diminuisce anziché crescere, perché l'elettrone da rimuovere non appartiene più a un orbitale $s$ ma a un orbitale $p$, meno penetrante e quindi meno legato. 
Analogamente, passando dal gruppo $5A$ al gruppo $6A$ si osserva una diminuzione, poiché il sesto elettrone deve appaiarsi in un orbitale $p$ già occupato e risente di una repulsione interelettronica aggiuntiva: la configurazione $np^3$, con i tre orbitali degeneri occupati da elettroni a spin parallelo secondo la regola di Hund, è particolarmente stabile.

Le energie di ionizzazione successive sono sempre crescenti, giacché ciascun elettrone viene strappato a uno ione di carica via via maggiore. 
L'aumento è graduale finché si rimuovono elettroni di valenza, mentre si osserva un salt di energia assai marcato nel momento in cui si comincia a intaccare un guscio interno completo. La posizione di tale salto individua il numero di elettroni di valenza dell'elemento, e costituisce una delle evidenze sperimentali più dirette dell'esistenza dei gusci.

\newpage

\section{Affinità elettronica}

Si definisce affinità elettronica\mkeyword{affinità elettronica} l'energia in gioco quando un atomo neutro isolato allo stato gassoso acquista un elettrone.

\begin{equation}
    X_{(g)} + e^- \longrightarrow X^-_{(g)}
    \label{eq:affinita}
\end{equation}

A differenza della \eqref{eq:ionizzazione}, il processo può essere sia esotermico sia endotermico. Negli alogeni, ai quali manca un solo elettrone per completare l'ottetto, esso è fortemente esotermico e l'affinità elettronica raggiunge i valori massimi in valore assoluto. 
Nei gas nobili, viceversa, l'elettrone dovrebbe collocarsi in un guscio nuovo e il processo risulta endotermico.

L'affinità elettronica cresce in generale procedendo verso destra lungo un periodo, con eccezioni in corrispondenza delle configurazioni particolarmente stabili $ns^2$, $ns^2np^3$ e $ns^2np^6$, per le quali l'elettrone aggiuntivo sarebbe fortemente destabilizzato.

\section{Metalli, non metalli e semimetalli}

Le proprietà finora esaminate permettono di classificare gli elementi in tre categorie, separate nella tavola periodica da una linea diagonale che scende dal boro verso l'astato.\\

I metalli\mkeyword{metalli} occupano la parte sinistra e centrale della tavola. Sono caratterizzati da energie di ionizzazione basse, e tendono pertanto a cedere elettroni formando cationi in soluzione acquosa. 
Allo stato solido riflettono la luce, sono duttili e malleabili e conducono bene la corrente elettrica e il calore, proprietà tutte riconducibili alla delocalizzazione degli elettroni di valenza nel reticolo. 
Il carattere metallico è massimo nei gruppi $1A$ e $2A$, con la sola eccezione dell'idrogeno, e cresce scendendo lungo un gruppo; tutti gli elementi di transizione sono metalli.\\

I non metalli\mkeyword{non metalli} occupano la parte destra della tavola.
Avendo elevata affinità elettronica, acquistano elettroni quando reagiscono con i metalli.

Allo stato solido sono fragili e non malleabili, e sono in genere cattivi conduttori. \\

I semimetalli\mkeyword{semimetalli}, o metalloidi, si collocano lungo la linea di separazione e presentano proprietà intermedie. 
Il silicio, ad esempio, ha aspetto metallico ma è fragile, e la sua conducibilità è intermedia fra quella dei metalli e quella degli isolanti: su questa proprietà si fonda l'intera elettronica a semiconduttori.

\subsection{Gruppo 1A: metalli alcalini}

Gli elementi del primo gruppo possiedono configurazione di valenza $ns^1$ e cedono con grande facilità l'unico elettrone esterno, essendo la loro energia di prima ionizzazione la più bassa dell'intero periodo. 
Nei composti compaiono pertanto sempre come cationi monovalenti $M^+$\mkeyword{metalli alcalini}. 
Allo stato elementare sono metalli morbidi e la loro reattività aumenta scendendo lungo il gruppo, poiché l'elettrone di valenza si trova progressivamente più lontano dal nucleo.\\

\subsection{Gruppo 2A: metalli alcalino-terrosi}

La configurazione di valenza $ns^2$ comporta la cessione di due elettroni e la formazione di cationi bivalenti $M^{2+}$\mkeyword{metalli alcalino-terrosi}. 
Rispetto agli alcalini essi sono più duri e meno reattivi, giacché la maggiore carica nucleare rende più difficoltosa la rimozione degli elettroni; anche in questo caso la reattività cresce scendendo lungo il gruppo.\\

\subsection{Idrogeno}

L'idrogeno occupa una posizione anomala nella tavola periodica. La configurazione $1s^1$ lo accomunerebbe ai metalli alcalini, ma il fatto che gli manchi un solo elettrone per completare il guscio lo avvicina agli alogeni, e in effetti ne condivide alcuni comportamenti.\mkeyword{idrogeno}\\

\subsection{Gruppo 6A: calcogeni}

Il sesto gruppo illustra con particolare chiarezza l'aumento del carattere metallico lungo un gruppo: l'ossigeno è un non metallo, il tellurio un semimetallo e il polonio un metallo a tutti gli effetti.\mkeyword{calcogeni}\\

L'ossigeno esiste in due forme allotropiche, la molecola biatomica $O_2$ e l'ozono $O_3$, quest'ultimo meno stabile e ottenibile dalla prima solo fornendo energia

\begin{equation}
    3O_{2(g)} \longrightarrow 2O_{3(g)} \qquad \Delta H = + 284,6 \ kJ
    \label{eq:ozono}
\end{equation}

condizione che si realizza nell'alta atmosfera per effetto della radiazione ultravioletta. Data l'elevata elettronegatività, l'ossigeno è un potente agente ossidante e presenta numero di ossidazione $-2$ nella gran parte dei composti, come nell'acqua, e $-1$ nei perossidi come $H_2O_2$. Lo zolfo, meno elettronegativo, compare nei composti binari come ione solfuro $S^{2-}$.

\subsection{Gruppo 7A: alogeni}

Gli alogeni possiedono configurazione $ns^2np^5$ e necessitano di un solo elettrone per completare l'ottetto: la loro affinità elettronica è pertanto la più elevata in valore assoluto\mkeyword{alogeni}. La reazione caratteristica è la riduzione a ione alogenuro

\begin{equation}
    X_2 + 2e^- \longrightarrow 2X^-
    \label{eq:alogenuro}
\end{equation}

Allo stato elementare formano molecole biatomiche $X_2$, il cui stato di aggregazione varia con la massa molecolare: fluoro e cloro sono gas, il bromo è liquido e lo iodio è solido, in conseguenza dell'aumento delle forze di London con il numero di elettroni. 
Il fluoro è fra le sostanze più reattive in assoluto, e la reattività diminuisce scendendo lungo il gruppo.

\subsection{Gruppo 8A: gas nobili}

I gas nobili possiedono i sottolivelli $s$ e $p$ completamente riempiti, configurazione che comporta la massima energia di ionizzazione e un'affinità elettronica praticamente nulla\mkeyword{gas nobili}. 
Non avendo alcuna tendenza né a cedere né ad acquistare elettroni, sono chimicamente inerti e si presentano monoatomici allo stato elementare, unico caso fra tutti gli elementi.\\

L'inerzia non è tuttavia assoluta: gli elementi più pesanti del gruppo, per i quali l'energia di ionizzazione è sensibilmente minore, formano composti con i soli elementi più elettronegativi, come nel caso di $XeF_4$ e $XeOF_4$.